\documentclass[12pt,a4paper]{article}
\usepackage[utf8]{inputenc}
\usepackage[T1]{fontenc}
\usepackage{geometry}
\usepackage{amsmath, amssymb}
\usepackage{booktabs}
\usepackage{listings}
\usepackage{xcolor}
\usepackage{hyperref}
\usepackage{parskip}

\geometry{margin=1in}

% ── Python Listing Style ──────────────────────────────────────────────────────
\definecolor{codebg}{RGB}{248,248,248}
\definecolor{codegreen}{rgb}{0,0.6,0}
\definecolor{codegray}{rgb}{0.5,0.5,0.5}
\definecolor{codeblue}{rgb}{0,0,0.8}
\definecolor{codeorange}{rgb}{0.8,0.4,0}

\lstdefinestyle{pythonstyle}{
    language=Python,
    backgroundcolor=\color{codebg},
    basicstyle=\ttfamily\footnotesize,
    keywordstyle=\color{codeblue}\bfseries,
    commentstyle=\color{codegreen}\itshape,
    stringstyle=\color{codeorange},
    numbers=left,
    numberstyle=\tiny\color{codegray},
    stepnumber=1,
    numbersep=8pt,
    tabsize=4,
    showstringspaces=false,
    breaklines=true,
    frame=single,
    captionpos=b,
    morekeywords={self, True, False, None, nn, F, torch, np, pd}
}
\lstset{style=pythonstyle}
% ─────────────────────────────────────────────────────────────────────────────

\title{\textbf{Dokumentasi Kode} \\ \large Perbandingan Strategi Portofolio Berbasis Graph dan GNN \\
\small \texttt{strategy\_comparison\_gnn\_thesis\_ready.ipynb}}
\author{Catatan Tesis}
\date{\today}

\begin{document}

\maketitle
\tableofcontents
\newpage

%=============================================================================
\section{Pendahuluan}
%=============================================================================
Notebook ini mengimplementasikan dan membandingkan sembilan strategi alokasi portofolio
pada data \textit{log-return} harian 9 aset kripto (periode 2017--2019). Pendekatan yang
dibandingkan mencakup:
\begin{enumerate}
    \item \textbf{Equally Weighted (EW)}: Bobot seragam $1/N$.
    \item \textbf{Classical Markowitz}: Optimasi \textit{mean-variance} standar.
    \item \textbf{Glasso Markowitz}: Markowitz dengan estimasi kovarians via \textit{Graphical Lasso}.
    \item \textbf{Network Markowitz} ($\gamma=0$ dan $\gamma=1$): Markowitz + sentralitas MST.
    \item \textbf{Graph Diversification} ($\theta=0.4$ dan $\theta=0.5$): Seleksi aset via MIS.
    \item \textbf{GNN Graph Diversification} (30 dan 50 epoch): MIS berbasis korelasi prediksi GNN.
\end{enumerate}

%=============================================================================
\section{Sel 1 -- Import Libraries}
%=============================================================================

\subsection{Deskripsi}
Mengimpor seluruh pustaka yang diperlukan. Pustaka utama yang digunakan:
\begin{itemize}
    \item \textbf{PyTorch \& PyTorch Geometric}: Arsitektur GCN dua lapis.
    \item \textbf{NetworkX}: Algoritma \textit{Maximum Independent Set} dan MST.
    \item \textbf{Scipy}: Optimasi \textit{quadratic programming} dan uji statistik.
    \item \textbf{Sklearn}: \textit{GraphicalLassoCV} untuk estimasi \textit{sparse} kovarians.
\end{itemize}

\subsection{Kode}
\begin{lstlisting}[caption={Import libraries}]
import numpy as np
import pandas as pd
import matplotlib.pyplot as plt
import matplotlib.gridspec as gridspec
import seaborn as sns
import networkx as nx
from scipy.optimize import minimize
from scipy.sparse.csgraph import minimum_spanning_tree
from scipy.linalg import eigh
from sklearn.covariance import GraphicalLassoCV
from scipy import stats
import torch
import torch.nn as nn
import torch.nn.functional as F
from torch_geometric.nn import GCNConv
from torch_geometric.data import Data
import warnings
warnings.filterwarnings('ignore')
plt.style.use('seaborn-v0_8-darkgrid')
sns.set_palette('husl')
print('Libraries imported successfully!')
print(f'PyTorch: {torch.__version__}')
\end{lstlisting}

%=============================================================================
\section{Sel 2 -- Load \& Prepare Real Crypto Data}
%=============================================================================

\subsection{Deskripsi}
Memuat data return harian dari file Excel \texttt{crypto\_data\_real.xlsx}.
Data diambil dari Yahoo Finance (Giudici et al., 2020) dan mencakup 9 aset kripto:
\textbf{BTC, ETH, XRP, BCH, LTC, BNB, EOS, XLM, TRX}.
USDT (stablecoin) dibuang karena return mendekati nol.
Baris dengan nilai nol (sebelum listing aset) juga dihapus.

Sub-periode pasar yang dianalisis:
\begin{itemize}
    \item \textbf{Bull (Nov--Des 2017)}: 52 hari.
    \item \textbf{Bear (2018)}: 364 hari.
    \item \textbf{Recovery (2019)}: 290 hari.
    \item \textbf{Full Period}: 706 hari (Nov 2017 -- Okt 2019).
\end{itemize}

\subsection{Kode}
\begin{lstlisting}[caption={Load dan preparasi data}]
EXCEL_FILE = 'crypto_data_real.xlsx'
df_raw = pd.read_excel(EXCEL_FILE, sheet_name='Returns', index_col=0)
df_raw.index = pd.to_datetime(df_raw.index)
df_raw.sort_index(inplace=True)
df_raw = df_raw.drop(columns=['USDT'], errors='ignore')  # stablecoin

# Hanya simpan baris di mana semua aset sudah aktif (return != 0)
mask = (df_raw != 0).all(axis=1)
df_returns = df_raw[mask].copy()

# Sub-period labels
BULL_END  = '2017-12-31'   # Bull: Nov 2017 - Des 2017
BEAR_END  = '2018-12-31'   # Bear: Jan 2018 - Des 2018
# Recovery: Jan 2019 - Okt 2019

periods = {
    'Bull (Nov-Dec 2017)' : df_returns.loc[df_returns.index <= BULL_END],
    'Bear (2018)'         : df_returns.loc[
                                (df_returns.index > BULL_END) &
                                (df_returns.index <= BEAR_END)],
    'Recovery (2019)'     : df_returns.loc[df_returns.index > BEAR_END],
    'Full Period'         : df_returns,
}

print('=== DATASET SUMMARY ===')
print(f'Assets  : {df_returns.columns.tolist()}')
print(f'Period  : {df_returns.index.min().date()} -> {df_returns.index.max().date()}')
print(f'Obs     : {len(df_returns)} trading days')
for pname, pdf in periods.items():
    print(f'{pname:<22}: {len(pdf):>4} days')
print(df_returns.describe().round(4))
\end{lstlisting}

%=============================================================================
\section{Sel 3 -- Helper Functions}
%=============================================================================

\subsection{Deskripsi}
Mendefinisikan fungsi-fungsi pembantu yang digunakan oleh semua strategi.

\subsection{3.1 \texttt{apply\_rmt\_filter} -- Random Matrix Theory Filter}
Memurnikan matriks korelasi empiris menggunakan teori matriks acak (RMT).
Nilai eigen di bawah batas Marchenko-Pastur $\lambda_+ = (1 + 1/\sqrt{Q})^2$
dianggap sebagai \textit{noise} dan di-nolkan.

\begin{lstlisting}[caption={\texttt{apply\_rmt\_filter}}]
def apply_rmt_filter(returns_data):
    T, N = returns_data.shape
    Q = T / N
    C = np.corrcoef(returns_data.T)
    eigenvalues, eigenvectors = eigh(C)
    # Balik urutan (terbesar ke terkecil)
    eigenvalues  = eigenvalues[::-1]
    eigenvectors = eigenvectors[:, ::-1]
    # Batas Marchenko-Pastur
    lambda_plus = 1 + (1/Q) + 2*np.sqrt(1/Q)
    # Nolkan nilai eigen yang dianggap noise
    Lambda_f = np.diag(np.where(eigenvalues > lambda_plus, eigenvalues, 0))
    return eigenvectors @ Lambda_f @ eigenvectors.T
\end{lstlisting}

\subsection{3.2 \texttt{build\_mst} -- Matriks Jarak MST}
Mengubah korelasi menjadi matriks jarak Euclidean untuk \textit{Minimum Spanning Tree}.
Formula jarak: $d_{ij} = \sqrt{2 - 2\rho_{ij}}$.

\begin{lstlisting}[caption={\texttt{build\_mst}}]
def build_mst(correlation_matrix):
    d = np.sqrt(2 - 2 * correlation_matrix)
    np.fill_diagonal(d, 0)
    return d
\end{lstlisting}

\subsection{3.3 \texttt{compute\_eigenvector\_centrality} -- Sentralitas Node}
Menghitung sentralitas \textit{eigenvector} dari graf jaringan. Aset dengan sentralitas
tinggi berada di pusat jaringan korelasi dan mendapat penalti lebih besar pada
Network Markowitz.

\begin{lstlisting}[caption={\texttt{compute\_eigenvector\_centrality}}]
def compute_eigenvector_centrality(distance_matrix):
    adj = 1 / (distance_matrix + 1e-8)   # adjacency = 1/jarak
    np.fill_diagonal(adj, 0)
    _, vecs = eigh(adj)
    pev = np.abs(vecs[:, -1])             # eigenvector terbesar
    return pev / pev.sum()                # normalisasi
\end{lstlisting}

\subsection{3.4 \texttt{get\_assets\_graph\_diversify} -- Seleksi MIS}
Membangun graf korelasi dan mengekstrak \textit{Maximum Independent Set} (set aset
yang tidak saling berkorelasi kuat). Edge ditambahkan jika $|\rho_{ij}| > \theta$.

\begin{lstlisting}[caption={\texttt{get\_assets\_graph\_diversify}}]
def get_assets_graph_diversify(returns_window, corr_threshold=0.4):
    corr_mat = returns_window.corr()
    G = nx.Graph()
    # Urutkan aset berdasarkan rata-rata return (terbesar duluan)
    assets = list(returns_window.mean().sort_values(ascending=False).index)
    G.add_nodes_from(assets)
    for i, a1 in enumerate(assets):
        for a2 in assets[i+1:]:
            if abs(corr_mat.loc[a1, a2]) > corr_threshold:
                G.add_edge(a1, a2)
    return list(nx.approximation.maximum_independent_set(G))
\end{lstlisting}

\subsection{3.5 Metrik Risiko}

\begin{lstlisting}[caption={Value at Risk, Rachev Ratio, Max Drawdown, Calmar Ratio}]
def calculate_var(returns, confidence=0.95):
    """Value at Risk pada level kepercayaan tertentu."""
    return np.percentile(returns, (1 - confidence) * 100)

def calculate_rachev_ratio(returns, alpha=0.10):
    """Rasio rata-rata tail atas terhadap tail bawah."""
    tu = np.percentile(returns, (1 - alpha) * 100)   # upper tail
    tl = np.percentile(returns, alpha * 100)           # lower tail
    cu = returns[returns >= tu].mean()
    cl = abs(returns[returns <= tl].mean())
    return cu / cl if cl > 0 else 0

def calculate_max_drawdown(cumulative_returns):
    """Drawdown maksimum dari nilai kumulatif portofolio."""
    running_max = np.maximum.accumulate(cumulative_returns)
    return ((cumulative_returns - running_max) / running_max).min()

def calculate_calmar_ratio(returns, cum_returns):
    """Calmar Ratio = Return Tahunan / |Max Drawdown|."""
    ann_ret = np.mean(returns) * 252
    mdd     = abs(calculate_max_drawdown(cum_returns))
    return ann_ret / mdd if mdd > 0 else 0
\end{lstlisting}

\subsection{3.6 \texttt{diebold\_mariano\_test} -- Uji DM}
Uji Diebold-Mariano (1995) untuk menguji signifikansi perbedaan dua strategi.
$H_0$: tidak ada perbedaan akurasi prediksi antara dua seri return.

\begin{lstlisting}[caption={\texttt{diebold\_mariano\_test}}]
def diebold_mariano_test(returns_a, returns_b, h=1):
    T   = len(returns_a)
    d   = returns_a - returns_b      # selisih return harian
    d_mean = np.mean(d)

    def autocov(xi, xm, k):
        T_ = len(xi)
        return np.sum((xi[:T_-k] - xm) * (xi[k:] - xm)) / T_ if T_ > k else 0

    var_d = autocov(d, d_mean, 0)
    for i in range(1, h):
        var_d += 2 * autocov(d, d_mean, i)

    dm = d_mean / np.sqrt(abs(var_d) / T) if var_d != 0 else 0
    pv = 2 * (1 - stats.norm.cdf(abs(dm)))
    return dm, pv
\end{lstlisting}

%=============================================================================
\section{Sel 4 -- Graph Neural Network (GCN Architecture)}
%=============================================================================

\subsection{Deskripsi}
Mendefinisikan arsitektur \textit{Two-layer Graph Convolutional Network} (GCN)
berdasarkan Kipf \& Welling (2017). Node features untuk setiap aset adalah
empat statistik dari return pada window saat itu:
\textbf{mean, std, skewness, kurtosis}.

\subsection{4.1 Kelas \texttt{CorrelationGNN}}
Model GCN dengan 2 lapis konvolusi. Output berupa \textit{embedding} tiap node
(aset) di ruang laten berdimensi \texttt{output\_dim=8}. Korelasi diprediksi
sebagai \textit{cosine similarity} antara pasangan \textit{embedding}.

\begin{lstlisting}[caption={Arsitektur CorrelationGNN}]
class CorrelationGNN(nn.Module):
    def __init__(self, input_dim=4, hidden_dim=16, output_dim=8):
        super().__init__()
        self.conv1 = GCNConv(input_dim, hidden_dim)
        self.conv2 = GCNConv(hidden_dim, output_dim)

    def forward(self, x, edge_index):
        x = F.relu(self.conv1(x, edge_index))
        x = F.dropout(x, p=0.2, training=self.training)
        x = F.relu(self.conv2(x, edge_index))
        return x

    def predict_correlation(self, emb):
        """Prediksi korelasi via cosine similarity embedding."""
        n = emb.shape[0]
        C = torch.zeros(n, n)
        for i in range(n):
            for j in range(i+1, n):
                s = F.cosine_similarity(emb[i].unsqueeze(0),
                                        emb[j].unsqueeze(0))
                C[i, j] = C[j, i] = s
        return C
\end{lstlisting}

\subsection{4.2 \texttt{create\_graph\_data} -- Pembuatan Data Graf}
Membangun objek PyG \texttt{Data} dari window return:
\begin{itemize}
    \item \textbf{Node features} $x$: [mean, std, skew, kurt] per aset.
    \item \textbf{Edges}: pasangan aset dengan $|\rho| > \text{corr\_threshold}$, ditambah arah sebaliknya (undirected).
    \item \textbf{Target labels} $y$: matriks korelasi empiris (digunakan sebagai target training).
\end{itemize}

\begin{lstlisting}[caption={\texttt{create\_graph\_data}}]
def create_graph_data(returns_window, corr_threshold=0.3):
    n = returns_window.shape[1]
    feats = np.stack([
        returns_window.mean().values,
        returns_window.std().values,
        returns_window.skew().values,
        returns_window.kurtosis().values
    ], axis=1)
    x    = torch.FloatTensor(feats)
    corr = returns_window.corr().values

    # Buat edge list dari pasangan yang berkorelasi
    edges = [[i, j] for i in range(n) for j in range(i+1, n)
             if abs(corr[i, j]) > corr_threshold]
    edges += [[j, i] for i, j in edges]   # undirected
    if not edges:                           # fallback: graf penuh
        edges = [[i, j] for i in range(n) for j in range(n) if i != j]

    ei = torch.LongTensor(edges).t().contiguous()
    return Data(x=x, edge_index=ei, y=torch.FloatTensor(corr))
\end{lstlisting}

\subsection{4.3 \texttt{train\_gnn} -- Pelatihan GNN}
Melatih GNN secara \textit{online} pada tiap window data menggunakan Adam optimizer
dan MSE loss antara korelasi yang diprediksi dan korelasi empiris.

\begin{lstlisting}[caption={\texttt{train\_gnn}}]
def train_gnn(model, data, epochs=50, lr=0.01):
    opt     = torch.optim.Adam(model.parameters(), lr=lr)
    loss_fn = nn.MSELoss()
    model.train()
    for _ in range(epochs):
        opt.zero_grad()
        emb  = model(data.x, data.edge_index)
        loss = loss_fn(model.predict_correlation(emb), data.y)
        loss.backward()
        opt.step()
    return model
\end{lstlisting}

%=============================================================================
\section{Sel 5 -- Portfolio Strategy Classes}
%=============================================================================

\subsection{5.1 \texttt{PortfolioStrategy} -- Kelas Dasar}
Kelas abstrak sebagai \textit{interface} semua strategi. Metode \texttt{get\_weights(r)}
menerima \texttt{DataFrame} return window dan mengembalikan array bobot portofolio.

\begin{lstlisting}[caption={Kelas dasar PortfolioStrategy}]
class PortfolioStrategy:
    def __init__(self, name):
        self.name = name
    def get_weights(self, r):
        raise NotImplementedError
\end{lstlisting}

\subsection{5.2 \texttt{EquallyWeighted}}
Bobot seragam $w_i = 1/N$ untuk semua $N$ aset.

\begin{lstlisting}[caption={EquallyWeighted}]
class EquallyWeighted(PortfolioStrategy):
    def get_weights(self, r):
        n = r.shape[1]
        return np.ones(n) / n
\end{lstlisting}

\subsection{5.3 \texttt{ClassicalMarkowitz}}
Optimasi mean-variance dengan constraint: $\sum w_i = 1$, $w_i \geq 0$,
dan ekspektasi return portofolio $\geq$ rata-rata return seluruh aset.
Minimisasi: $\min_w \; w^\top \Sigma w$.

\begin{lstlisting}[caption={ClassicalMarkowitz}]
class ClassicalMarkowitz(PortfolioStrategy):
    def get_weights(self, r):
        n  = r.shape[1]
        mu = r.mean().values
        S  = r.cov().values
        f  = lambda w: w @ S @ w
        c  = [{'type': 'eq',  'fun': lambda w: np.sum(w) - 1},
              {'type': 'ineq','fun': lambda w: w @ mu - mu.mean()}]
        res = minimize(f, np.ones(n)/n, method='SLSQP',
                       bounds=[(0,1)]*n, constraints=c)
        return res.x if res.success else np.ones(n)/n
\end{lstlisting}

\subsection{5.4 \texttt{GlassoMarkowitz}}
Sama dengan Markowitz klasik, namun matriks kovarians diestimasi menggunakan
\textit{Graphical Lasso} yang menghasilkan matriks \textit{sparse} (sehingga
merefleksikan struktur dependensi yang lebih bersih).

\begin{lstlisting}[caption={GlassoMarkowitz}]
class GlassoMarkowitz(PortfolioStrategy):
    def __init__(self, name='Glasso Markowitz', alpha=0.01):
        super().__init__(name)
        self.alpha = alpha

    def get_weights(self, r):
        n  = r.shape[1]
        mu = r.mean().values
        try:
            g = GraphicalLassoCV(alphas=[self.alpha], cv=3)
            g.fit(r.values)
            S = g.covariance_
        except:
            S = r.cov().values   # fallback ke kovarians biasa
        f  = lambda w: w @ S @ w
        c  = [{'type': 'eq',  'fun': lambda w: np.sum(w) - 1},
              {'type': 'ineq','fun': lambda w: w @ mu - mu.mean()}]
        res = minimize(f, np.ones(n)/n, method='SLSQP',
                       bounds=[(0,1)]*n, constraints=c)
        return res.x if res.success else np.ones(n)/n
\end{lstlisting}

\subsection{5.5 \texttt{NetworkMarkowitz}}
Mengintegrasikan informasi \textit{network} (sentralitas eigenvector) ke dalam
fungsi objektif Markowitz sebagai penalti:
\[
\min_w \; w^\top \Sigma_f w + \gamma \sum_i c_i w_i
\]
di mana $\Sigma_f$ adalah kovarians berbasis RMT-filtered correlation,
dan $c_i$ adalah sentralitas eigenvector aset $i$.

\begin{lstlisting}[caption={NetworkMarkowitz}]
class NetworkMarkowitz(PortfolioStrategy):
    def __init__(self, name='Network Markowitz', gamma=0):
        super().__init__(name)
        self.gamma = gamma

    def get_weights(self, r):
        n    = r.shape[1]
        mu   = r.mean().values
        sig  = r.std().values
        Cf   = apply_rmt_filter(r)
        dm   = build_mst(Cf)
        cent = compute_eigenvector_centrality(dm)
        Sf   = np.outer(sig, sig) * Cf   # Kovarians dari korelasi RMT
        f    = lambda w: w @ Sf @ w + self.gamma * np.sum(cent * w)
        c    = [{'type': 'eq',  'fun': lambda w: np.sum(w) - 1},
                {'type': 'ineq','fun': lambda w: w @ mu - mu.mean()}]
        res  = minimize(f, np.ones(n)/n, method='SLSQP',
                        bounds=[(0,1)]*n, constraints=c)
        return res.x if res.success else np.ones(n)/n
\end{lstlisting}

\subsection{5.6 \texttt{GraphDiversification}}
Memilih subset aset yang saling \textit{independent} (tidak berkorelasi tinggi)
menggunakan \textit{Maximum Independent Set}. Bobot didistribusikan secara merata
di antara aset terpilih.

\begin{lstlisting}[caption={GraphDiversification}]
class GraphDiversification(PortfolioStrategy):
    def __init__(self, name='Graph Diversification', corr_threshold=0.4):
        super().__init__(name)
        self.corr_threshold = corr_threshold

    def get_weights(self, r):
        n      = r.shape[1]
        assets = r.columns.tolist()
        sel    = get_assets_graph_diversify(r, self.corr_threshold)
        w      = np.zeros(n)
        if sel:
            for a in sel:
                w[assets.index(a)] = 1.0 / len(sel)
        else:
            w = np.ones(n) / n
        return w
\end{lstlisting}

\subsection{5.7 \texttt{GNNGraphDiversification}}
Versi GD yang menggunakan korelasi hasil prediksi GCN (bukan korelasi empiris)
untuk membangun graf dan menentukan \textit{threshold} secara adaptif:
\[
\theta^* = \text{clip}\bigl(\bar{|\rho_{GNN}|} + 0.5 \cdot \sigma_{|\rho_{GNN}|},\ 0.3,\ 0.7\bigr)
\]

\begin{lstlisting}[caption={GNNGraphDiversification}]
class GNNGraphDiversification(PortfolioStrategy):
    def __init__(self, name='GNN GD', base_threshold=0.4, train_epochs=30):
        super().__init__(name)
        self.base_threshold = base_threshold
        self.train_epochs   = train_epochs
        self.gnn_model      = CorrelationGNN(input_dim=4, hidden_dim=16, output_dim=8)

    def get_weights(self, r):
        n      = r.shape[1]
        assets = r.columns.tolist()

        # 1) Buat graph data & latih GNN
        gd = create_graph_data(r, corr_threshold=0.3)
        self.gnn_model = train_gnn(self.gnn_model, gd,
                                   epochs=self.train_epochs)
        # 2) Prediksi korelasi dari embedding GCN
        self.gnn_model.eval()
        with torch.no_grad():
            emb = self.gnn_model(gd.x, gd.edge_index)
            pc  = self.gnn_model.predict_correlation(emb)

        pn = pc.numpy()
        up = pn[np.triu_indices(n, k=1)]

        # 3) Threshold adaptif
        thr = np.clip(np.mean(np.abs(up)) + 0.5*np.std(np.abs(up)), 0.3, 0.7)

        # 4) Bangun graf dengan threshold adaptif
        G       = nx.Graph()
        asorted = list(r.mean().sort_values(ascending=False).index)
        G.add_nodes_from(asorted)
        for i, a1 in enumerate(asorted):
            for a2 in asorted[i+1:]:
                i1, i2 = assets.index(a1), assets.index(a2)
                if abs(pn[i1, i2]) > thr:
                    G.add_edge(a1, a2)

        # 5) MIS & equal weighting
        sel = list(nx.approximation.maximum_independent_set(G))
        w   = np.zeros(n)
        if sel:
            for a in sel:
                w[assets.index(a)] = 1.0 / len(sel)
        else:
            w = np.ones(n) / n
        return w
\end{lstlisting}

%=============================================================================
\section{Sel 6 -- Backtesting Framework}
%=============================================================================

\subsection{Deskripsi}
Fungsi umum untuk menjalankan \textit{rolling-window backtest}. Parameter kunci:

\begin{table}[h!]
\centering
\begin{tabular}{ll}
\toprule
Parameter & Nilai \\
\midrule
\texttt{window\_size} & 120 hari (training window) \\
\texttt{rebalance\_freq} & 7 hari (mingguan) \\
\texttt{transaction\_cost} & 0.1\% per rebalancing \\
\bottomrule
\end{tabular}
\caption{Parameter backtesting}
\end{table}

\subsection{6.1 \texttt{backtest\_strategy}}

\begin{lstlisting}[caption={\texttt{backtest\_strategy}}]
def backtest_strategy(strategy, df_ret, window_size=120,
                      rebalance_freq=7, transaction_cost=0.001):
    port_returns = []
    dates        = []
    for i in range(window_size, len(df_ret), rebalance_freq):
        train = df_ret.iloc[i-window_size:i]
        w     = strategy.get_weights(train)
        test_end = min(i + rebalance_freq, len(df_ret))
        test     = df_ret.iloc[i:test_end]
        for j in range(len(test)):
            dr = np.dot(w, test.iloc[j].values)
            if j == 0 and port_returns:
                dr -= transaction_cost   # biaya transaksi saat rebalancing
            port_returns.append(dr)
            dates.append(test.index[j])

    res = pd.DataFrame({'date': dates, 'return': port_returns})
    res['cumulative_return'] = (1 + res['return']).cumprod()
    return {
        'strategy'          : strategy.name,
        'returns'           : np.array(port_returns),
        'cumulative_returns': res['cumulative_return'].values,
        'dates'             : dates,
        'results_df'        : res
    }
\end{lstlisting}

\subsection{6.2 \texttt{calculate\_metrics}}
Menghitung seluruh metrik kinerja dari hasil backtest.

\begin{lstlisting}[caption={\texttt{calculate\_metrics}}]
def calculate_metrics(result):
    r  = result['returns']
    cr = result['cumulative_returns']
    ann_r = np.mean(r) * 252              # Return tahunan
    ann_v = np.std(r) * np.sqrt(252)      # Volatilitas tahunan
    sharpe = ann_r / ann_v if ann_v > 0 else 0
    return {
        'Strategy'          : result['strategy'],
        'Total Return (%)'  : round((cr[-1] - 1) * 100, 2),
        'Annual Return (%)' : round(ann_r * 100, 2),
        'Annual Vol (%)'    : round(ann_v * 100, 2),
        'Sharpe Ratio'      : round(sharpe, 4),
        'VaR 95% (%)'       : round(calculate_var(r, 0.95) * 100, 4),
        'Rachev Ratio'      : round(calculate_rachev_ratio(r, 0.10), 4),
        'Max Drawdown (%)'  : round(calculate_max_drawdown(cr) * 100, 2),
        'Calmar Ratio'      : round(calculate_calmar_ratio(r, cr), 4),
    }
\end{lstlisting}

%=============================================================================
\section{Sel 7 -- Menjalankan Backtest (Full Period)}
%=============================================================================

\subsection{Deskripsi}
Mendefinisikan semua instance strategi dan menjalankan backtest secara berurutan
pada keseluruhan data. Hasilnya disimpan dalam dictionary \texttt{results}.

\begin{lstlisting}[caption={Menjalankan backtest untuk semua strategi}]
strategies = [
    EquallyWeighted('Equally Weighted'),
    ClassicalMarkowitz('Classical Markowitz'),
    GlassoMarkowitz('Glasso Markowitz', alpha=0.01),
    NetworkMarkowitz('Network Markowitz (gamma=0)',   gamma=0),
    NetworkMarkowitz('Network Markowitz (gamma=1.0)', gamma=1.0),
    GraphDiversification('Graph Divers. (theta=0.4)', corr_threshold=0.4),
    GraphDiversification('Graph Divers. (theta=0.5)', corr_threshold=0.5),
    GNNGraphDiversification('GNN Graph Divers. (30 ep)', train_epochs=30),
    GNNGraphDiversification('GNN Graph Divers. (50 ep)', train_epochs=50),
]

print('Running backtests on real crypto data...')
results = {}
for s in strategies:
    print(f'  Testing: {s.name}...')
    results[s.name] = backtest_strategy(
        s, df_returns, window_size=120, rebalance_freq=7)
    print(f'  Done')
print('\nAll backtests completed.')
\end{lstlisting}

%=============================================================================
\section{Sel 8 -- Performance Metrics}
%=============================================================================

\subsection{Deskripsi}
Menghitung dan menampilkan tabel metrik performa untuk semua strategi. Kolom
terbaik per metrik di-\textit{highlight} dengan warna hijau.

\begin{lstlisting}[caption={Tampilkan tabel metrik}]
metrics_df = pd.DataFrame(
    [calculate_metrics(r) for r in results.values()]
).set_index('Strategy')

def highlight_best(col):
    """Highlight nilai terbaik per kolom."""
    is_min = col.name in ['Annual Vol (%)', 'VaR 95% (%)', 'Max Drawdown (%)']
    best   = col.min() if is_min else col.max()
    return ['background-color: #d4edda; font-weight: bold'
            if v == best else '' for v in col]

print('PERFORMANCE METRICS (full period)')
print('='*100)
print(metrics_df.to_string())
print('='*100)
metrics_df.style.apply(highlight_best)
\end{lstlisting}

%=============================================================================
\section{Sel 9 -- Validasi Statistik (Matriks DM)}
%=============================================================================

\subsection{Deskripsi}
Membangun matriks $N \times N$ pairwise DM-test untuk semua pasangan strategi.
Nilai DM positif berarti strategi pada baris mengalahkan strategi pada kolom.
$p < 0.05$ menandakan perbedaan yang signifikan secara statistik.

\begin{lstlisting}[caption={Komputasi matriks DM test}]
strategy_names = list(results.keys())
n_strat = len(strategy_names)

dm_stat_mat = pd.DataFrame(index=strategy_names,
                            columns=strategy_names, dtype=float)
dm_pval_mat = pd.DataFrame(index=strategy_names,
                            columns=strategy_names, dtype=float)

for i, s1 in enumerate(strategy_names):
    for j, s2 in enumerate(strategy_names):
        if i == j:
            dm_stat_mat.loc[s1, s2] = 0.0
            dm_pval_mat.loc[s1, s2] = 1.0
        else:
            dm, pv = diebold_mariano_test(
                results[s1]['returns'], results[s2]['returns'])
            dm_stat_mat.loc[s1, s2] = round(dm, 3)
            dm_pval_mat.loc[s1, s2] = round(pv, 4)

sig_matrix = dm_pval_mat < 0.05   # True = signifikan
print('DM Test P-Value Matrix (p < 0.05 = significant)')
print(dm_pval_mat.to_string())
print()
print('DM Test Statistic Matrix (positive = row outperforms column)')
print(dm_stat_mat.to_string())
\end{lstlisting}

%=============================================================================
\section{Sel 10 -- Heatmap DM Test}
%=============================================================================

\subsection{Deskripsi}
Visualisasi matriks DM sebagai \textit{heatmap} dengan warna:
\begin{itemize}
    \item \textbf{Hijau}: statistik DM positif (baris mengalahkan kolom).
    \item \textbf{Merah}: statistik DM negatif (baris kalah dari kolom).
    \item \textbf{Bintang (*)} pada sel: $p < 0.05$ (signifikan).
\end{itemize}

\begin{lstlisting}[caption={Heatmap DM test}]
# Singkat nama strategi agar muat di heatmap
short_names = {
    'Equally Weighted'              : 'EW',
    'Classical Markowitz'           : 'CMV',
    'Glasso Markowitz'              : 'Glasso',
    'Network Markowitz (gamma=0)'   : 'Net-MV(0)',
    'Network Markowitz (gamma=1.0)' : 'Net-MV(1)',
    'Graph Divers. (theta=0.4)'     : 'GD(0.4)',
    'Graph Divers. (theta=0.5)'     : 'GD(0.5)',
    'GNN Graph Divers. (30 ep)'     : 'GNN-30',
    'GNN Graph Divers. (50 ep)'     : 'GNN-50',
}
dm_stat_r = dm_stat_mat.rename(index=short_names, columns=short_names).astype(float)
pval_r    = dm_pval_mat.rename(index=short_names, columns=short_names).astype(float)

fig, ax = plt.subplots(figsize=(11, 8))
mask_diag = np.eye(len(dm_stat_r), dtype=bool)
sns.heatmap(dm_stat_r, annot=True, fmt='.2f', center=0, cmap='RdYlGn',
            linewidths=0.5, ax=ax, mask=mask_diag,
            cbar_kws={'label': 'DM Statistic (positive = row outperforms col)'})

# Tambah tanda '*' untuk sel yang signifikan
for i in range(len(dm_stat_r)):
    for j in range(len(dm_stat_r.columns)):
        if i != j and pval_r.iloc[i, j] < 0.05:
            ax.text(j+0.85, i+0.15, '*', color='black',
                    fontsize=14, fontweight='bold')

ax.set_title('Diebold-Mariano Test Statistic Matrix\n'
             '(* = significant at p<0.05, green = row outperforms)',
             fontsize=12)
ax.set_xlabel('Benchmark Strategy')
ax.set_ylabel('Test Strategy')
plt.tight_layout()
plt.savefig('dm_test_heatmap.png', dpi=150, bbox_inches='tight')
plt.show()
print('Saved: dm_test_heatmap.png')
\end{lstlisting}

%=============================================================================
\section{Sel 11 -- Visualisasi Cumulative Returns}
%=============================================================================

\subsection{Deskripsi}
Plot perbandingan nilai kumulatif portofolio setiap strategi selama periode penuh
beserta \textit{shading} untuk pembagian fase pasar (Bull/Bear/Recovery).
File disimpan ke \texttt{cumulative\_returns\_comparison.png}.

\begin{lstlisting}[caption={Plot cumulative returns (ringkasan logika)}]
fig, ax = plt.subplots(figsize=(14, 7))

# Plot setiap strategi
for name, res in results.items():
    ax.plot(res['dates'], res['cumulative_returns'], label=name)

# Tambah shading fase pasar
ax.axvspan(df_returns.index.min(), pd.Timestamp(BULL_END),
           alpha=0.1, color='green', label='Bull (2017)')
ax.axvspan(pd.Timestamp(BULL_END), pd.Timestamp(BEAR_END),
           alpha=0.1, color='red', label='Bear (2018)')
ax.axvspan(pd.Timestamp(BEAR_END), df_returns.index.max(),
           alpha=0.1, color='blue', label='Recovery (2019)')

ax.axhline(y=1, color='black', linestyle='--', alpha=0.5)
ax.set_title('Cumulative Returns -- All Strategies (Real Crypto Data)')
ax.set_xlabel('Date')
ax.set_ylabel('Portfolio Value (start = 1.0)')
ax.legend(loc='upper right', fontsize=8)
plt.tight_layout()
plt.savefig('cumulative_returns_comparison.png', dpi=150, bbox_inches='tight')
plt.show()
\end{lstlisting}

%=============================================================================
\section{Ringkasan Arsitektur Sistem}
%=============================================================================

\begin{center}
\begin{tabular}{lll}
\toprule
\textbf{Komponen} & \textbf{Teknologi} & \textbf{Fungsi} \\
\midrule
Data & Yahoo Finance (Excel) & 9 aset kripto, log-return harian \\
Optimasi Portofolio & SciPy SLSQP & Mean-variance optimization \\
Korelasi Sparse & sklearn GraphicalLasso & Kovarians lebih bersih \\
Network Analysis & NetworkX & MST, MIS, eigenvector centrality \\
RMT Filter & SciPy eigh & Pembersihan noise dari kovarians \\
GNN & PyTorch Geometric GCNConv & Prediksi korelasi adaptif \\
Backtesting & Pandas / NumPy & Rolling-window simulation \\
Validasi & SciPy stats (DM-test) & Signifikansi perbedaan strategi \\
\bottomrule
\end{tabular}
\end{center}

\end{document}
